针对 Chrono 的非光滑接触（non-smooth contact, NSC）框架在高精度连续区域接触建模中的局限，本文提出一种面向 SDF 几何表征的多 patch / 多 subpatch 保持 wrench 等效的接触约化方法。该方法以符号距离场为几何基础，首先构造高密度接触区域参考模型；随后依据连通性、法向变化、几何曲率与 wrench 误差将接触区域自适应分解为多个 patch 与 subpatch；最后针对每个 subpatch 构建少量代表点组成的 reduced support set，使其在给定误差阈值下逼近 dense 接触模型的可行 wrench 锥，并将这些代表点重新注入 Chrono NSC 的标准点接触求解链路中。与传统基于少量缓存点的接触流形维护方式不同，本文方法不是简单进行几何删点，而是将接触约化重新表述为“由 dense admissible wrench cone 到 reduced admissible wrench cone 的误差控制投影”。在理论上，本文给出了局部共面 subpatch 的三点 wrench 完备性、法向可行性与中心压力点支撑域之间的几何条件，并提出一套以 wrench 误差、CoP 误差、内部 gap 误差和拓扑误差共同驱动的自适应细化策略。本文进一步给出面向凸轮连续区域接触、点接触以及长条形接触（齿轮、插销--轴承等）的统一实验设计框架。该方法的目标并非逐点重建真实压力场，而是在保持 Chrono NSC 全局互补求解框架不变的前提下，提高连续接触区域在刚体动力学层面的合力、合力矩与作用线等效精度，为高保真 SDF 接触与现有多体动力学求解器的耦合提供一条可落地的技术路线。
